\section{GoalScope}
\label{app:goalscope}
\subsection{Capture positions}
\arnab{this should be absorbed in evaluation}
\label{app:pivots}
The evaluation indexes the final four prompt tokens around a prefilled frontier assignment and the first generated implementation token: previous-line end, indentation, variable, assignment operator, and first implementation. Position names refer to the actual tokenized boundary; they are not interchangeable with an entire source line. The first implementation token is not a pre-implementation pivot. A cached state at a given position has no causal access to subsequent tokens, even when those tokens are displayed for context.


% \begin{figure}[ht]\centering
% \includegraphics[width=.92\linewidth]{Figures/author/evaluation_setup.pdf}
% \caption{\textbf{Capture positions and the two probes, enlarged.} Pink marks the previous-line end, purple indentation, yellow the variable, light green the assignment operator, and dark green the first implementation token. The orange and green arrows illustrate faithfulness and causality readouts from the variable and assignment positions, respectively. In the evaluation, each probe is run separately at all five positions. Only source activations up to the selected pivot enter its cached vector.}
% \label{fig:pivot-probes-detail}
% \end{figure}

\subsection{Prompts for Faithfulness and Causality}
\arnab{full prompts for faithfulness and causality}

\subsection{Steering setup}
\arnab{discuss design choices: replace vs additive\_norm. steering layers for the additive\_norm strategy, lambda for both. we didn't do a systematic search over the hparams.}
At source position $t$, the main readout aggregates projected head contributions $g_t=\sum_{(\ell,j)\in H}a_t^{\ell j}$. At probe layer $\ell$ and current position $s$, norm-relative additive injection uses
\begin{equation}
r_{\ell,s}\leftarrow r_{\ell,s}+\alpha\|r_{\ell,s}\|\frac{g_t}{\|g_t\|}.
\end{equation}
The main configuration uses $\alpha=.1$, layers 35--45 inclusive, and repeated injection during greedy probe generation. This coefficient is a fraction of the probe's local residual norm and is not interchangeable with an unnormalized additive coefficient. The source vector is fixed for the whole readout.

The default verbalization prompt is:
\begin{quote}\small
System: You are a helpful assistant proficient in programming.\\
User: What are you doing at this line of the code. Keep your answer as brief as possible.\\
Assistant:
\end{quote}
\arnab{full faithfulness and causality prompts. use consistent box format throughout}
The probe receives no source code as text. Implementation transfer uses a code-stub probe instead and measures nontrivial agreement with the source subgoal. A distinct replacement variant substitutes selected heads' pre-output-projection activations in their corresponding layers, rather than aggregating them into a single direction; this is the variant used in the cross-model table. Its layer sites come from the head set, not the additive injection band.

\subsection{Evaluation}
\arnab{a subpart of the evaluation setup. with the full prompt in \Cref{app:oracle_prompts}}
An oracle judges whether the generated description or implementation matches a nontrivial part of the extracted frontier code. The full system rubrics, user-message templates, and saved examples are in \Cref{app:oracle_prompts}. Reported implementation-transfer results apply trivial-readout sanitization to avoid crediting a probe's default sign/zero comparisons. These judgments do not certify execution correctness. Clean continuation agreement, probe readout agreement, and intervention-induced changes are separate measurements. Corpus size and capture-position pooling are stated with each result.
